\documentclass[11pt]{article}

\usepackage[utf8]{inputenc}
\usepackage[T1]{fontenc}
\usepackage[spanish,es-tabla]{babel}
\usepackage{lmodern}
\usepackage[a4paper,margin=2cm]{geometry}
\usepackage{array}
\usepackage{booktabs}
\usepackage{xcolor}
\usepackage{enumitem}
\usepackage{hyperref}

\setlength{\parskip}{0.5em}
\setlength{\parindent}{0pt}
\renewcommand{\arraystretch}{1.3}
\setlist{nosep,leftmargin=1.4em}

\definecolor{marca}{HTML}{0B3A5C}
\hypersetup{colorlinks=true,linkcolor=marca,urlcolor=marca}

% Enlaces al prototipo publicado en GitHub Pages.
\newcommand{\urlapp}{https://thankstochatgpt.github.io/LOARD}
\newcommand{\app}[1]{\href{\urlapp#1}{\texttt{#1}}}

\title{Cobertura de historias de usuario\\[0.3em]
\large Loard - Hito 2 - Módulo I: Prototipo de alta fidelidad}
\author{}
\date{17 de noviembre de 2026}

\begin{document}
\maketitle

\section{Introducción}

Este documento muestra en qué parte del prototipo se puede ver cada una de las 12 historias de usuario del MVP. El prototipo está publicado en \href{\urlapp/}{\nolinkurl{https://thankstochatgpt.github.io/LOARD/}} y cada ruta de la tabla es un enlace directo.

Para entrar se elige un perfil de ejemplo; el correo y la contraseña (\texttt{Demo2026!}) se completan solos. Los datos son ficticios y se guardan en el navegador.

\section{Matriz de historias y pantallas}

\begin{center}
\small
\begin{tabular}{>{\raggedright\arraybackslash}p{1.4cm} >{\raggedright\arraybackslash}p{4.4cm} >{\raggedright\arraybackslash}p{6.4cm} >{\raggedright\arraybackslash}p{3cm}}
\toprule
\textbf{HU} & \textbf{Historia} & \textbf{Dónde se ve} & \textbf{Responsable} \\ \midrule
HU-M01 & Iniciar sesión según el rol & \app{/acceso} & Juan David Peña \\
HU-M02 & Registrar el corte de obra & \app{/cortes}, \app{/cortes/nuevo} & Juan David Peña \\
HU-M03 & Adjuntar fotos y soportes al corte & \app{/cortes/ejemplo-3} (Corte \#08) & Sebastián Gutiérrez \\
HU-M04 & Visar o devolver el corte & \app{/cortes} con el perfil Director & Sebastián Gutiérrez \\
HU-M05 & Publicar documentos y armar la checklist & \app{/dossier}, \app{/dossier/documentos} & Angel Inyena Pasuy \\
HU-M06 & Cargar soportes y buscar documentos & \app{/dossier}, \app{/dossier/documentos} & Angel Inyena Pasuy \\
HU-M07 & Registrar ensayos NDT & \app{/dossier/ensayos} & Juan David Peña \\
HU-M08 & Avance y entrega del dossier & \app{/dossier/avance} & Juan David Peña \\
HU-M09 & Cursos y aptitud médica (SST) & \app{/sst} & Jonathan López \\
HU-M10 & Personal del corte con bloqueo SST & \app{/cortes/nuevo/personal} & Jonathan López \\
HU-M11 & Bitácora de cambios & \app{/dossier/bitacora} & Ronald Chávez \\
HU-M12 & Tablero del proyecto & \app{/dashboard} & Ronald Chávez \\
\bottomrule
\end{tabular}
\end{center}

Las 12 historias tienen al menos una pantalla en el prototipo.

\section{Qué se puede ver en cada historia}

Para cada historia se resume qué hace el usuario y cómo responde el prototipo cuando algo sale mal (un dato incompleto, un archivo no permitido o la falta de señal en la obra).

\subsection*{HU-M01: Iniciar sesión según el rol}
Al elegir un perfil de ejemplo se llenan solos el correo y la contraseña. Cada perfil ve en el menú solo los módulos de su trabajo.
\begin{itemize}
  \item Si la contraseña está mal, aparece un mensaje claro sin decir cuál de los dos datos falló.
  \item Si alguien abre por dirección una pantalla que no le corresponde, ve un aviso de acceso restringido y un enlace para volver al inicio.
  \item El botón con el símbolo de apagar cierra la sesión.
\end{itemize}

\subsection*{HU-M02: Registrar el corte de obra}
El residente registra el frente, el periodo, las actividades con su cantidad, el personal con sus horas y los insumos usados.
\begin{itemize}
  \item Si falta un dato, el formulario dice exactamente qué completar antes de guardar.
  \item Si sale sin guardar, al volver puede continuar o descartar el borrador.
  \item Cada corrección pide un motivo y crea una versión nueva; el historial conserva quién cambió qué y cuándo.
  \item Un corte ya visado no se puede editar y la pantalla explica por qué.
\end{itemize}

\subsection*{HU-M03: Adjuntar fotografías y soportes al corte}
El residente agrega fotos y PDF al corte, con el botón o arrastrándolos, y los abre en un visor sin salir del corte.
\begin{itemize}
  \item Un archivo de otro tipo o de más de 15 MB se rechaza con la causa y cómo solucionarlo.
  \item Si se cae la señal, el archivo queda en espera y se sube solo al volver la conexión.
  \item Un soporte quitado por error se recupera con ``Deshacer'' durante 5 segundos.
  \item El Director ve los mismos soportes en modo consulta.
\end{itemize}

\subsection*{HU-M04: Visar o devolver el corte}
El Director entra a una bandeja con los cortes pendientes, revisa cantidades, horas, insumos y soportes en una sola pantalla, y decide.
\begin{itemize}
  \item Visar pide confirmación; después el corte ya no se puede modificar.
  \item Devolver exige escribir una observación (hay causas frecuentes para elegir rápido). El residente la ve, corrige y reenvía.
  \item Sin señal, los botones de decisión quedan bloqueados con el motivo visible.
  \item Tener cortes pendientes nunca impide que el residente registre cortes nuevos.
\end{itemize}

\subsection*{HU-M05: Publicar documentos vigentes y armar la checklist}
Calidad publica los formatos y procedimientos del sistema de gestión y arma la lista de requisitos de las seis carpetas del dossier.
\begin{itemize}
  \item Publicar un documento con un código existente crea una versión nueva; la anterior queda en el historial.
  \item Retirar un requisito pide confirmación y explica que deja de contar en el avance.
\end{itemize}

\subsection*{HU-M06: Cargar soportes y buscar documentos}
A cada requisito se le carga su PDF, que Calidad valida u observa con un motivo. Los documentos se buscan por nombre, código, carpeta o estado.
\begin{itemize}
  \item Una validación hecha por error se puede deshacer durante unos segundos.
  \item Si un PDF no se carga, el mensaje invita a reintentar o a elegir otro archivo.
\end{itemize}

\subsection*{HU-M07: Registrar ensayos NDT ligados a una junta}
Calidad registra el ensayo de una junta soldada: laboratorio, número de informe, método, fecha, resultado y el PDF del laboratorio.
\begin{itemize}
  \item Cada informe crea su propio requisito en la carpeta 4 del dossier y no reemplaza el de otra junta.
  \item No se permiten informes repetidos ni fechas posteriores al día de la demostración.
  \item Junto al resultado se aclara que ``Conforme'' no significa que el documento ya esté validado por Calidad.
\end{itemize}

\subsection*{HU-M08: Monitorear el avance del dossier y exportarlo}
Una pantalla muestra el porcentaje validado, el avance por carpeta, lo que falta y los días que quedan del plazo de diez días desde la aceptación del acta.
\begin{itemize}
  \item Solo cuentan los requisitos validados; un archivo cargado todavía no suma.
  \item Siempre se puede descargar un borrador parcial en ZIP; la entrega final se habilita cuando todo está validado.
\end{itemize}

\subsection*{HU-M09: Registrar cursos y aptitud médica con semáforo}
Coordinación SST ve la lista de trabajadores con los casos críticos primero, y registra cursos y conceptos de aptitud médica.
\begin{itemize}
  \item Un semáforo avisa a los 30 y a los 15 días de vencer, y bloquea al trabajador cuando vence.
  \item Al elegir las fechas de un curso se ve cómo quedará el semáforo antes de guardar.
  \item La aptitud médica es solo Habilitado o No Habilitado, sin datos clínicos. Los cursos no se borran: se corrigen registrando uno nuevo.
\end{itemize}

\subsection*{HU-M10: Asignar personal al corte con bloqueo SST}
Al agregar personal a un corte, el residente ve a todos los trabajadores: los disponibles arriba y los bloqueados abajo, con candado y el motivo.
\begin{itemize}
  \item Un trabajador bloqueado no se puede marcar; si se intenta, la fila resalta el motivo.
  \item Un resumen muestra los trabajadores elegidos con sus alertas de vencimiento.
  \item Lo que se cambia en SST se refleja de inmediato en este selector.
\end{itemize}

\subsection*{HU-M11: Consultar la bitácora de cambios}
Calidad elige un proyecto y un corte o requisito, y ve todos sus cambios de estado: quién, cuándo, de qué estado a cuál y por qué, del más reciente al más antiguo.
\begin{itemize}
  \item Es solo de consulta: ningún evento se puede editar ni borrar.
  \item Si no hay eventos o la consulta falla, el prototipo lo dice y permite reintentar.
\end{itemize}

\subsection*{HU-M12: Consultar el tablero de avance}
Director y Gerencia ven el avance físico del proyecto frente al tiempo del contrato, el porcentaje validado del dossier, los días que quedan del plazo y las seis carpetas.
\begin{itemize}
  \item Cada carpeta abre el detalle de sus requisitos, pero desde el tablero no se puede modificar nada.
\end{itemize}

\section{Conclusión}

El prototipo cubre todas las historias del MVP con pantallas que comparten el mismo diseño y los mismos datos: el estado SST se refleja en los cortes, y el dossier alimenta la bitácora y el tablero. Los detalles de diseño están en los documentos del Módulo II y las pruebas con usuarios en el \emph{Reporte de validación RITE} del Módulo III.

\end{document}
